\documentclass[11pt,a4paper]{article}
\usepackage[utf8]{inputenc}
\usepackage[T1]{fontenc}
\usepackage[french]{babel}
\usepackage{geometry}
\usepackage{graphicx}
\usepackage{xcolor}
\usepackage{booktabs}
\usepackage{array}
\usepackage{float}
\usepackage{hyperref}
\geometry{margin=2.2cm}
\hypersetup{colorlinks=true,linkcolor=blue,urlcolor=blue}
\setlength{\parskip}{5pt}
\setlength{\parindent}{0pt}

\newcommand{\capture}[2]{%
  \begin{figure}[H]\centering
  \includegraphics[width=0.96\linewidth,height=0.48\textheight,keepaspectratio]{#2}
  \caption{#1}
  \end{figure}
}

\title{\textbf{Laboratoire 5}\\Segmentation et partitionnement des traces GPS}
\author{Étudiant : \textbf{MOUAD MARZAQ}\\
Professeure : \textbf{GHYZLANE CHERRADI}}
\date{10 novembre 2026}

\begin{document}
\maketitle
\tableofcontents
\newpage

\section{Objectif}
Ce laboratoire compare plusieurs façons d'organiser et d'interroger un historique
de positions GPS sous PostgreSQL/PostGIS : indexation spatiale et temporelle,
segmentation logique en bandes, partitions journalières, tuilage et réorganisation
physique. Deux requêtes servent de fil conducteur : \textbf{Q1}, qui compte les
positions du 1\textsuperscript{er} novembre situées à moins de 500 m de la
sous-station 1, et \textbf{Q2}, qui recherche les véhicules à moins de 2 km d'un
point d'incident sur la ligne haute tension 1 pendant la dernière heure du 2 novembre.

Les géométries sont exprimées en EPSG:32629, un système métrique. Les positions
du jeu d'essai sont synthétiques : 200 véhicules, un relevé toutes les deux
minutes et sept jours, soit 1\,008\,000 lignes. Les performances dépendent du
matériel, du cache, de la version PostgreSQL/PostGIS et des données réelles.

\section{Préparation et protocole}
Le script \texttt{lab5.sql} crée un schéma isolé nommé
\texttt{lab5\_soumission}; il ne supprime pas les tables sources
\texttt{public.sous\_stations} et \texttt{public.lignes\_haute\_tension}.
L'exécution se fait avec \texttt{psql} sur la base de données du laboratoire.
La commande utilisée doit afficher les ordres et les temps :

\begin{verbatim}
psql -h localhost -p 5433 -U postgres `
  -d ReseauElectriqueSIG -v ON_ERROR_STOP=1 `
  -a -f lab5.sql
\end{verbatim}

Pour les comparaisons de plans, le script répète les mesures : la seconde
exécution (cache plus chaud) est celle à reporter. Les temps ne sont pas
comparables entre machines; les plans \texttt{EXPLAIN (ANALYZE, BUFFERS)}
permettent aussi de relever le type de parcours et les pages lues ou trouvées
dans le cache. Les valeurs mesurées et les captures doivent provenir de
l'exécution réelle; les tableaux ci-dessous reprennent les résultats lisibles
dans les captures fournies, et signalent les mesures absentes.

\section{Task 1 — Référence et index}
\subsection{Jeu de données}
La fonction de génération place une majorité de points dans un voisinage
aléatoire des sous-stations et le reste dans l'emprise du réseau. Le contrôle
attendu est de 1\,008\,000 positions, du 27 octobre à 00:00 au 2 novembre à
23:58. Les index GiST sur la géométrie et B-tree sur la date fournissent deux
chemins d'accès adaptés aux prédicats spatiaux et temporels.

\capture{Contrôle du volume et de la période du jeu de données}{captures/shot_233526.png}

\subsection{Mesures Q1 et Q2}
\begin{table}[H]
\centering
\begin{tabular}{lrrrr}
\toprule
Requête & Lignes & Temps (ms) & Buffers hit & Buffers read \\
\midrule
Q1, table de référence & 6\,331 & 357,539 & 2\,802 & 7\,879 \\
\bottomrule
\end{tabular}
\caption{Q1 de référence selon le plan \texttt{EXPLAIN (ANALYZE, BUFFERS)} capturé.}
\end{table}

\capture{Plan d'exécution de Q1 sur la table non partitionnée}{captures/shot_234040.png}
\capture{Taille de la table et des index après création}{captures/shot_233929.png}

\textbf{Interprétation.} Un index spatial permet de retrouver rapidement les
identifiants de lignes proches, mais il n'assure pas que les lignes de la table
soient rangées spatialement sur disque. Un parcours peut donc consulter de
nombreuses pages du tas pour récupérer les géométries et les autres colonnes.
L'index ne supprime pas non plus le coût de la vérification exacte
\texttt{ST\_DWithin}. Le plan et les buffers observés permettent de distinguer
les pages réellement lues des pages déjà présentes en mémoire.

\section{Task 2 — Segmentation logique en trois zones}
La coordonnée Y de l'emprise est divisée en trois intervalles de même hauteur :
Sud (1), Centre (2) et Nord (3). La colonne \texttt{zone\_id} matérialise cette
segmentation; l'index composite \texttt{(zone\_id, date)} vise les requêtes qui
combinent zone et période.

Les résultats capturés confirment que les 1\,008\,000 positions ont toutes
reçu une zone (aucune valeur NULL). La sous-station 1 se trouve dans la zone 2.
Après l'ajout de \texttt{zone\_id}, la taille mesurée passe de 81 à 171 Mo pour
la table et de 47 à 100 Mo pour les index; le \texttt{VACUUM} n'a pas réduit
ces valeurs dans la mesure disponible.

\capture{Taille après ajout de la colonne de zone}{captures/shot_234604.png}
\capture{Vérification du nombre de zones NULL}{captures/shot_234919.png}
\capture{Zone attribuée à la sous-station 1}{captures/shot_234938.png}

Le filtre \texttt{zone\_id = 2} de Q1 ne conserve que 3\,571 positions,
contre 6\,331 pour la requête de référence. La segmentation en bandes n'est
donc pas une restriction exacte du disque de 500 m : des points du voisinage
de la station se trouvent de l'autre côté d'une frontière de zone. Ce filtre
ne doit pas être utilisé seul pour répondre à une requête spatiale exacte.
Le plan capturé sans index composite dure 413,981 ms et lit 10\,736 blocs
depuis le stockage; le nombre de lignes retournées confirme l'écart de
résultat.

\capture{Plan de Q1 avec filtre logique sur la zone}{captures/shot_234315.png}

\section{Task 3 — Partitionnement journalier}
La table \texttt{gps\_part} est partitionnée par plages de dates, avec une
partition par jour et une partition par défaut. Un prédicat sur \texttt{date}
permet l'élagage des partitions avant leur lecture; sans ce prédicat, une
condition uniquement spatiale ne permet généralement pas de choisir un jour et
les partitions restantes doivent être considérées.

\begin{table}[H]
\centering
\begin{tabular}{lrrrr}
\toprule
Q1 & Lignes & Temps (ms) & Buffers hit & Buffers read \\
\midrule
Table de référence & 6\,331 & 357,539 & 2\,802 & 7\,879 \\
Partition journalière & 6\,331 & 91,193 & 1\,481 & 60 \\
\bottomrule
\end{tabular}
\caption{Comparaison des plans Q1 capturés. Les deux plans retournent le même compte.}
\end{table}

L'insertion dans une table partitionnée vise directement une seule partition.
La taille des index à maintenir et le volume de pages concernés peuvent donc
être réduits. Le résultat réel dépend toutefois des index, des écritures WAL,
du cache et du système de stockage. Pour la rétention, supprimer une partition
ancienne par \texttt{DROP TABLE} est habituellement plus direct qu'effacer ses
lignes une à une par \texttt{DELETE}; une purge par \texttt{DELETE} laisse aussi
des espaces morts jusqu'à la maintenance et ne garantit pas une réduction
immédiate du fichier.

\capture{Effectifs des partitions journalières}{captures/shot_000424.png}
\capture{Plan Q1 sur la partition du 1\textsuperscript{er} novembre}{captures/shot_000707.png}
\capture{Plan avec borne temporelle dynamique}{captures/shot_000547.png}

La requête dynamique est exécutée avant l'ajout du 3 novembre dans le script :
elle se termine donc à la dernière position alors disponible, le 2 novembre à
23:58. Elle illustre l'élagage à l'exécution, sans correspondre exactement à la
fenêtre Q2 de 23:00 à minuit. Les captures disponibles ne montrent pas les
durées comparatives d'insertion et de purge; elles ne sont donc pas chiffrées ici.

\section{Task 4 — Tuilage et GeoHash}
\subsection{Grille de 2 km}
La grille carrée sert à agréger les positions par cellule et à construire une
couche de densité. Une jointure \texttt{ST\_Intersects} peut attribuer une
position située exactement sur une frontière à plusieurs cellules. Le calcul
par \texttt{floor(ST\_X/2000)} et \texttt{floor(ST\_Y/2000)} donne au contraire
une affectation unique, à condition d'utiliser le même origine de grille.

\begin{table}[H]
\centering
\begin{tabular}{lrrrr}
\toprule
Méthode & Cellules non vides & Min. & Max. & Moyenne / écart-type \\
\midrule
Grille 2 km & 20 & 11 & 30\,922 & 7\,200,00 / 9\,156,57 \\
\bottomrule
\end{tabular}
\caption{Densité du 1\textsuperscript{er} novembre.}
\end{table}

\capture{Plan de calcul des densités par cellules}{captures/shot_001355.png}
\capture{Statistiques des cellules non vides}{captures/shot_001422.png}

\subsection{Filtre GeoHash de Q2}
Le GeoHash est calculé après transformation en EPSG:4326. Un préfiltre par
préfixe réduit les candidats, mais une cellule GeoHash n'est pas un cercle de
rayon 2 km. Le prédicat \texttt{ST\_DWithin} doit donc rester dans la requête
pour confirmer la distance exacte. Le rapport de différences entre Q2 de
référence et Q2 optimisée sert à détecter les faux négatifs introduits par un
ensemble de cellules candidates trop étroit.

Les quatre coins d'un rectangle suffisent pour retrouver les bornes min/max
d'une grille carrée alignée, mais cela ne signifie pas que les quatre GeoHash
des coins énumèrent automatiquement toutes les cellules GeoHash touchées. La
couverture doit être testée (ou calculée en ajoutant les cellules voisines)
avant d'utiliser ce préfiltre en production. Le tuilage complète un index GiST;
il ne le remplace pas pour une recherche de voisinage exacte.

\capture{Plan Q2 avec préfiltre GeoHash et distance exacte}{captures/shot_002005.png}

Dans cette exécution, le plan GeoHash renvoie 200 véhicules distincts en
1\,430,833 ms. Le préfiltre est une approximation : \texttt{ST\_DWithin} reste
indispensable pour confirmer la distance. Les performances affichées ne sont
pas directement comparables à Q2 de référence, car les captures ne documentent
pas une répétition contrôlée des mêmes conditions de cache.

\section{Task 5 — Clustering physique}
\texttt{CLUSTER} réécrit une table en suivant l'ordre d'un index existant.
L'expérience utilise un index \texttt{(vehicule\_id, date)} pour le 30 octobre,
un index GeoHash pour le 31 octobre et un index GiST pour le 1\textsuperscript{er}
novembre. La corrélation de \texttt{pg\_stats} indique si l'ordre physique des
lignes ressemble à l'ordre logique d'une colonne. Une corrélation proche de
\(1\) signale une organisation croissante; une valeur proche de \(-1\) indique
l'ordre inverse.

\begin{table}[H]
\centering
\begin{tabular}{llr}
\toprule
Partition & Ordre de clustering & Corrélation après \\
\midrule
30 octobre & \texttt{vehicule\_id} & 1,000 \\
31 octobre & \texttt{geohash} & 1,000 \\
1\textsuperscript{er} novembre & GiST sur la géométrie; date & 0,004 \\
\bottomrule
\end{tabular}
\caption{Valeurs de corrélation après \texttt{CLUSTER}, selon \texttt{pg\_stats}.}
\end{table}

Le clustering est une réorganisation ponctuelle : les insertions suivantes
peuvent dégrader progressivement l'ordre. Son intérêt doit être confirmé par les
plans Q1/Q2 et le coût de maintenance.

\capture{Corrélations mesurées après clustering}{captures/shot_002819.png}

\section{Task 6 — Synthèse et choix d'architecture}
\textbf{Conclusion.} Pour un historique qui croît continuellement et dont les
requêtes portent sur une période, le partitionnement par date constitue une
base adaptée : il facilite l'élagage temporel et la purge par période. Chaque
partition peut conserver un index GiST pour les recherches spatiales et un
index temporel/composite selon les requêtes mesurées. Un index GeoHash peut
compléter cette organisation si le préfiltre est assez sélectif et si sa
couverture est validée contre Q2. Le conserver sur chaque ligne augmente le
volume de stockage et le coût d'écriture; il faut donc comparer ce coût au gain
mesuré. Enfin, un clustering périodique n'est à retenir que si les tests
montrent un bénéfice durable malgré les nouvelles insertions.

Les valeurs de temps et de buffers correspondent à l'environnement de
l'exécution capturée et peuvent varier avec le matériel, le cache et la version
de PostgreSQL/PostGIS. Les mesures d'insertion et de purge n'étant pas visibles
dans les captures fournies, elles restent à relever si une comparaison chiffrée
de ces opérations est exigée.

\end{document}
